% A neutrino corridor for a 10 TeV muon collider at Fermilab
% PRL-format concept letter. Compile: pdflatex (revtex4-2) x2 + bibtex if split out.
\documentclass[aps,prl,reprint,superscriptaddress,nofootinbib]{revtex4-2}
\usepackage{graphicx}
\usepackage{amsmath}
\usepackage[colorlinks=true,linkcolor=blue,citecolor=blue,urlcolor=blue]{hyperref}

\begin{document}

\title{A Neutrino Corridor for a 10~TeV Muon Collider at Fermilab}

\author{L.~Lee}
\affiliation{University of Tennessee, Knoxville, TN, USA}

\date{\today}

\begin{abstract}
The decay neutrinos of a multi-TeV muon collider form pencil beams of
unprecedented intensity and energy, and their interaction with the
surrounding landscape --- where they graze the Earth's surface --- is a
recognized siting constraint. We invert the constraint into an asset. We
show that an 11~km collider ring fits the Fermilab campus with its
interaction-point straight on the true-north meridian $-88.222972^\circ$,
tilted upward by 15.4~mrad (one LEP of ring-plane inclination). The
north-going neutrino beam then leaves the ground \emph{inside the
laboratory fence}, delivering ${\sim}3\times10^{10}$ TeV-scale neutrino
interactions per tonne-year to an on-site experiment on the present
IMCC-class optics --- whose beam divergence, not the decay kinematics,
sets the spot size~\cite{MINT} --- and ${\sim}2.6\times10^{12}$ if a
dedicated low-divergence drift is built, before climbing a utility
right-of-way into federally navigable airspace; the south-going
beam dives beneath the region's aquifers to a perigee of 789~m and
resurfaces 198~km away on University of Illinois agricultural land, where
its exit strip doubles as a second neutrino experiment of
${\sim}10^{6}$--$10^{8}$ interactions per tonne-year. The second interaction point's
south beam exits onto the same university's airport. Every point at which
a beam passes within 500~ft of the ground surface lies on federal,
utility-easement, or university land, and every potable-aquifer crossing
beyond the laboratory fence is bounded at a factor ${\gtrsim}5\times10^{3}$
(${\gtrsim}45$ even with the dedicated drift) below drinking-water
limits. The
alignment costs ${\sim}11\%$ of one accelerator ring's circumference and
one LEP-scale plane tilt, and converts the muon collider's defining
radiological liability into two shared-science facilities.
\end{abstract}

\maketitle

\emph{Introduction.}---A muon collider at
$\sqrt{s}=10$~TeV is the most compact route to the
energy frontier now under serious study~\cite{MCForum,IMCC}. Its defining
environmental feature is also its most distinctive experimental
opportunity: every muon decays, and the decays in each long straight
section emit neutrinos along the straight's axis --- a cone of half-angle
$1/\gamma \simeq 21~\mu$rad (at $E_\mu = 5$~TeV) for a FODO straight,
broadened to an effective $\sigma_\theta \simeq 0.15$~mrad for the
interaction-point straight by the final-focus optics~\cite{MINT}.
Where that plume intersects the Earth's surface tens of kilometres
away, the showers from neutrino interactions in the final metres of soil
can produce annual doses approaching regulatory limits --- the
``neutrino radiation'' problem first quantified by King~\cite{King} and
adopted as a driving constraint by the International Muon Collider
Collaboration, whose reference siting aims its insertion plumes at the
Mediterranean and a fenced mountainside~\cite{IMCC}. The same beam,
sampled \emph{before} it exits, is the most intense high-energy neutrino
beam ever conceived, with fully known flavor composition and parent
kinematics~\cite{NuSlice}.

This Letter reports a complete siting geometry for the Fermilab campus in
which both properties are engineered simultaneously and deliberately: the
beam's two ground intersections are \emph{chosen}, both fall on
institutional land, both host experiments, and the intervening 200~km of
private land is crossed only at altitudes or depths at which no surface
owner's use of their property is touched.

\emph{The geodetic coincidence.}---Three surveyed facts make the geometry
work. (i)~A commercial transmission right-of-way (ComEd
\emph{Aurora--Wayne}) runs almost exactly due north--south at longitude
$\simeq -88.223^\circ$ from the Fermilab boundary northward.
(ii)~The northeast corner of the laboratory's site polygon lies on that
meridian to within one metre of longitude. (iii)~Along the meridian
$-88.222972^\circ$ the site offers 5.48~km of north--south extent ---
room for an 11~km racetrack's 700~m straights and 1.5~km-radius arcs, and
for the acceleration chain's straights on the same line. Because a
due-north geodesic holds constant longitude, aligning every long straight
with the corridor reduces to placing it on this meridian at azimuth
$0^\circ$.

\emph{Machine geometry.}---We fit an 11.0~km racetrack collider
($R_{\rm arc}=1528$~m, $2\times700$~m insertions) against the real site
boundary with a 150~m setback, its east straight on the corridor meridian
and its interaction point (IP) at $41.84428^\circ$\,N, placed such that
both ends of the straight and both beam emergences remain well inside the
boundary. The ring plane is tilted $\theta_0 = 15.40$~mrad
($0.88^\circ$) upward toward north --- 10\% beyond LEP's built precedent
of 14~mrad --- with the straightaway at 191~m above sea level, 34~m below
local grade, at the glacial-drift/Silurian-dolomite contact. The rapid-cycling-synchrotron
chain (6.28 and 14.72~km racetracks at 60--80~m depth) shares the
meridian with mrad-scale plane pitches so that all straight plumes
converge on one on-site detector volume; the alignment's only machine
cost is a $10.8\%$ circumference shortfall in the largest RCS relative to
an unconstrained boundary-hugging fit.

In a spherical-geoid frame a straight beam launched at height $h_0$ with
elevation angle $\theta_0$ has height
$h(s)=h_0+\theta_0 s+s^{2}/2R_\oplus$ above the curving surface. Northward,
with real (SRTM) terrain, the beam surfaces at $41.8640^\circ$\,N ---
2.2~km from the IP and 655~m \emph{inside} the laboratory fence --- crosses
the boundary 8~m above grade, climbs the right-of-way, passes the nearest
residential area 116~m (380~ft) above grade --- 350~ft above a 10~m
roofline --- and enters federally navigable airspace (500~ft above ground)
9.3~km beyond the fence.
Southward the same straight aims $15.40$~mrad below the horizon: the beam
passes beneath the Aurora conurbation at 74--152~m depth, reaches a
perigee of 789~m in the brackish upper Mt.~Simon sandstone --- the horizon
northern Illinois uses for natural-gas storage, sealed beneath the Eau
Claire aquitard and below every potable unit --- and re-emerges
$2\theta_0$ of central angle later: 198~km south, at $40.067^\circ$\,N,
on the University of Illinois Urbana--Champaign (UIUC) \emph{South Farms},
approaching under the city of Urbana at 66--119~m depth. A
$\pm0.1$~mrad trim moves this exit $\pm1.3$~km along the meridian:
targeting a chosen parcel is a survey problem, not a physics one. The
diametrically opposite second IP, on the ring's west straight (meridian
$-88.2598^\circ$), inherits the same plane tilt; its south beam exits at
${\simeq}40.045^\circ$\,N on Willard Airport --- owned and operated by the
same university --- after a $+0.13$~mrad trim well inside the
$\pm1$~mrad mover budget already carried by the IMCC design~\cite{IMCC}.

\emph{The civic envelope.}---The siting rule that organizes the study is
geometric: wherever either beam is within 500~ft (152~m) of the ground
surface, vertically, the surface land must belong to an institutional
partner --- the Department of Energy, the utility easement, or the
university. Above 500~ft over uncongested land an aircraft may legally
fly (14~CFR~91.119); since \emph{United States v.\ Causby} (1946) a
landowner's airspace claim reaches only the ``immediate reaches'' they
can use, so a silent, invisible, depositing-nothing neutrino pencil
hundreds of feet up is comfortably in the public highway. Below 500~ft of
rock the beam is deeper than every private well (the shallow Silurian
aquifer's wells run 30--90~m), and --- unlike the deep tunnels bored
under Chicago through thousands of recorded subsurface
easements --- nothing is built: neutrinos simply pass. Auditing the
baseline against this rule leaves four honest exceptions, all deep or
high crossings, none touching the surface: 1.9~km of overflight at
384--500~ft over low-density countryside; 5.6~km at 74--152~m under
Aurora's east edge; and the farmland-plus-Urbana approach at 66--152~m.
Standard avigation or subsurface easements close each, though none
plausibly requires it.

\emph{Fluxes and rates.}---With IMCC-class beam parameters
($1.8\times10^{12}$ muons per bunch per sign at 5~Hz, $1.2\times10^{7}$~s
operational year, 90\% transmission, and 85\% of stored muons decaying
within each 0.2~s store), $5.3\times10^{18}$ decays per year are aimed
along each direction of the IP straight ($700$~m of an 11~km ring). The
pencil-limit on-axis flux is $\Phi = N\gamma^{2}/\pi L^{2}$, but on the
present optics the plume is divergence-smeared: MINT~\cite{MINT}, decaying
muons along the IMCC v0.6+v0.9 interaction-region lattice, finds
$\sigma_\theta \simeq 0.15$~mrad $\simeq 7/\gamma$, diluting the on-axis
density by $2\gamma^{2}\sigma_\theta^{2} \simeq 10^{2}$ and washing out
the energy--angle correlation, so any point in the core sees the
angle-integrated spectrum, $\langle E_\nu\rangle \simeq 0.33\,E_\mu$.
An on-site detector 1.3~km from the IP then intercepts
${\sim}3\times10^{10}$ charged- and neutral-current interactions per
tonne-year on axis (CSMS cross-sections, both species per decay) --- six
orders of magnitude beyond present TeV-neutrino samples, in a beam whose
flavor content ($\nu_\mu\bar\nu_e$ or $\bar\nu_\mu\nu_e$, selectable by
circulation direction) is exactly known --- and
${\sim}2.6\times10^{12}$ if a fraction $f \simeq 0.5$ of the straight is
rebuilt as a dedicated dispersion-free drift (a waisted
$\beta^{*}\!\sim\!\ell$ optic with $\sigma_\theta = \sqrt{\epsilon/\ell}
\simeq 1~\mu$rad), the one machine-design ask this concept adds. At the
South Farms exit, 198~km out, the surviving flux delivers
${\sim}1.3\times10^{6}$ ($f=0$) to ${\sim}10^{8}$ ($f\simeq0.5$)
interactions per tonne-year: the exit strip that must be controlled
anyway becomes a far detector site with $L/E$ two orders of magnitude
beyond any accelerator baseline, on land owned by a research university.
The aligned RCS chain adds a time-gated 63~GeV--5~TeV energy chirp
carrying ${\sim}40\%$ of the collider's rate at the deep-hall option,
since its FODO straights emit true $1/\gamma$ pencils~\cite{site}.

\emph{Beam size.}---Muon decay is isotropic in the muon rest frame, so a
single decaying muon emits the beamed profile
$dN/d\Omega = N\gamma^{2}/[\pi(1+\gamma^{2}\theta^{2})^{2}]$ about its own
direction. Along the IP straight those directions are not the axis: a
field-free drift conserves angles, the final-focus drifts are
${\sim}150$~m long at $\sigma_\theta = \sqrt{\epsilon_N/\gamma\beta^{*}}
= 0.59$~mrad ($\epsilon_N = 25~\mu$m, $\beta^{*} = 1.5$~mm), and the
matching sections contribute dispersively, so the full-lattice simulation
of Ref.~\cite{MINT} finds an effective plume divergence of
0.1--0.2~mrad~$\gg 1/\gamma$ with the energy--angle correlation (the
``prism'') washed out. This study adopts $\sigma_\theta = 0.15$~mrad: the
50\% radius is then 24~cm at the on-site detector, 52~cm at the fence,
and 36~m at the southern exit (the kinematic 2.7~cm / 4.2~m survive only
for the dedicated-drift pencil fraction $f$). A one-tonne tungsten
cylinder no longer swallows the beam --- at $f=0$ it intercepts 8\% ---
so the on-site detector is better modelled on the low-density
tracking-target benchmarks of Ref.~\cite{MINT}, sized to the core with a
front monitor for the ${\sim}2$ rock muons per bunch crossing that
reference computes. At the exits the cross-section is stretched along the
ground by $1/\sin(\mathrm{graze})$: the northern emergence paints a
0.8~m~$\times$~50~m streak inside the laboratory fence, while the
southern one is a 72~m~$\times$~4.6~km strip, which is the actual
geometry of the land-use question --- wider than the pencil model gave,
with its peak dose diluted by the same factor of ${\sim}10^{2}$.

\emph{Radiological audit.}---Dose: everywhere along the corridor except
the two chosen exits the beam is deep underground or high overhead, and
the exit strips themselves are the entire land-use footprint; at the
southern exit the $1/L^{2}$ dilution over 198~km and the measured plume
divergence bring even King's deliberately conservative equilibrium
model~\cite{King} to ${\sim}0.8$~mSv/yr with no further mitigation
(${\sim}0.2$~mSv/yr with a $\pm0.5$~mrad vertical segmentation) on fenced
farmland; the dedicated-drift pencil would concentrate ${\sim}40$~mSv/yr
into a 10~m core there, making detector rate and exit dose a single
design knob (the framework, corrections, and the deep-reference
alternative are in the accompanying study~\cite{site}). Groundwater: neutrinos
activate nothing directly; the only pathway is the hadronic showers of
the ${\sim}2\times10^{15}$ interactions per year distributed along the
full 198~km chord ($3.5\times10^{-4}$ interaction lengths of rock),
yielding ${\sim}6\times10^{16}$ tritium atoms per year
(${\simeq}0.05$~Ci/yr at saturation) with a $1/L^{2}$ weighting toward
the IP. In the maximally conservative limit --- groundwater stagnant
inside the decimetres-wide beam core at each aquifer's closest approach,
for many tritium lifetimes, with zero exchange --- the EPA drinking-water
limit (740~Bq/L) is exceeded only within ${\sim}0.2$~km of the IP
(${\sim}1.4$~km in the dedicated-drift scenario), entirely
inside the laboratory boundary. Traced through a dipping stratigraphic
model of the Illinois Basin's flank (unit tops anchored to NuMI
construction geology, ISGS mapping, and the region's gas-storage fields),
the municipal sandstone aquifers are crossed ${\sim}14$~km and
${\sim}28$~km out and ceiling at ${\sim}5\times10^{3}$ and
${\sim}2\times10^{4}$ below the limit ($45\times$ and $190\times$ with
the dedicated drift); the chord bottoms in the brackish Mt.~Simon; and
the drift sands near Champaign--Urbana ceiling at ${\sim}10^{6}\times$
below.
Groundwater advection --- metres per day in the karstic on-site
dolomite --- lowers all of these by further orders of magnitude.

\emph{Conclusion.}---The neutrino radiation problem of a 10~TeV muon
collider at Fermilab can be solved by construction rather than
mitigation: one LEP of ring-plane tilt and a meridian alignment place
both surface intersections of the collision straight --- and both of the
second interaction point's --- on federal or university land, thread the
200~km between them outside the envelope of private use, bound aquifer
activation orders of magnitude below drinking-water standards, and
produce two neutrino experiments of unprecedented reach as a by-product.
The geometry uses only surveyed data and analytic conservative models; a
real lattice, FLUKA/MARS transport, and hydrogeological modelling are the
obvious next steps, and all inputs, code, and an interactive
reproduction of every number in this Letter are public~\cite{site}.

\begin{acknowledgments}
Geometry, terrain, and boundary data: OpenStreetMap contributors (ODbL)
and NASA SRTM. This concept study used no beam time and harmed no
groundwater.
\end{acknowledgments}

\begin{thebibliography}{9}

\bibitem{MCForum}
K.~M. Black \emph{et al.}, ``Muon Collider Forum report,''
JINST \textbf{19}, T02015 (2024), arXiv:2209.01318.

\bibitem{IMCC}
C.~Accettura \emph{et al.} (International Muon Collider Collaboration),
``Interim report for the International Muon Collider Collaboration,''
CERN Yellow Rep. Monogr. (2024), arXiv:2407.12450.

\bibitem{King}
B.~J. King, ``Potential hazards from neutrino radiation at muon
colliders,'' arXiv:physics/9908017.

\bibitem{MINT}
J.-Y. Choi, M. Hostert, P. Li, and Z. Liu,
``The forward neutrino flux and its secondaries at a 10~TeV muon
collider,'' arXiv:2608.02718.

\bibitem{NuSlice}
L.~Bojorquez-Lopez, M.~Hostert, T.~Menzo, and J.~Zupan,
``The neutrino slice at muon colliders,'' arXiv:2412.14115.

\bibitem{site}
L.~Lee, ``FNAL muon collider corridor siting study,''
\url{https://lawrenceleejr.github.io/FNALMuonColliderSite/} and
\url{https://github.com/lawrenceleejr/FNALMuonColliderSite} (2026).

\end{thebibliography}

\end{document}
